\chapter{Representability and Internalization}

\section{Numeral Encoding}

To internalize metatheoretic statements into the formal system of Peano Arithmetic, we must first be able to refer to external natural numbers within the syntax. This is achieved by mapping each metatheoretic natural number to a canonical closed term.

\begin{defn}\label{defn:numerals}
    The \textsl{numeral encoding} is a function $\fun{num}\colon\N\to\fun{Term}$, typically denoted by an overline $n\mapsto\num{n}$, defined by structural recursion on $\N$:
    \begin{align*}
        \fun{num}(0)&\defeq 0\\
        \fun{num}(\fun{succ}(n))&\defeq s(\fun{num}(n)).
    \end{align*}
    Thus, $\num n$ is a syntax tree consisting of $n$ applications of the successor constructor~$s$ to the constant constructor~$0$.
\end{defn}

\begin{prop}\label{prop:num-gives-closed-terms}
 For every natural number\/ $n\colon\N$, the numeral\/ $\num{n}$ is a closed term.
\end{prop}

\begin{proof}
    We proceed by induction on $n\colon\N$. For the base case $n\equiv0$, the numeral is $\num{0}\equiv 0\colon\fun{Term}$. By Definition~\ref{defn:FV}, $\fun{FV}(0)(y)\defeq\fun0$, which implies that $\neg\fun{FV}(0)(y)$ is inhabited for any variable $y\colon\fun{Var}$.
    
    For any $y\colon\fun{Var}$, we have the following definitional equalities:
    \begin{align*}
        \fun{FV}(\num{\fun{succ}(n)})(y)
            &\equiv\fun{FV}(s(\num{n}))(y)\\
            &\equiv\fun{FV}(\num{n})(y).
    \end{align*}
    Then
    \[
        \fun{isClosed}(\num{\fun{succ}(n)})
            \equiv\fun{isClosed}(\num n),
    \]
    which is inhabited by the inductive hypothesis.
\end{proof}

\begin{lem}\label{lem:num-is-embedding}
 The numeral encoding function is an embedding. Consequently, for any\/ $n,m\colon\N$, we have an equivalence:
 \[
 (\num{n}\eq{\fun{Term}}\num{m}) \simeq (n\eq{\N}m).
 \]
\end{lem}

\begin{proof}
    We proceed by double induction on $n,m\colon\N$.
    
    When $n\equiv0$ and $m\equiv0$, we have the constant map $x\mapsto\fun{refl}_0$. Since $\N$ is also a set, it suffices to construct a map $(\num{n}\eq{\fun{Term}}\num{m})\to(n\eq{\N}m)$. By Theorem~\ref{thm:syntax-are-sets}, $\fun{Term}$ is a set. Since $\N$ is also a set, the map is an equivalence.
    
    When $n\equiv0$ and $m\equiv\fun{succ}(m')$, the domain equality type becomes
    \[
        0\eq{\fun{Term}}s(\num{m'}).
    \]
    By Theorem~\ref{thm:constructor-disjointness} this type is inhabited. Since $0\eq\N\fun{succ}(m')$ is also inhabited, both types are equivalent.

    The symmetric case follows similarly.
    
    When $n\equiv\fun{succ}(n')$ and $m\equiv\fun{succ}(m')$, Theorem~\ref{thm:constructor-injectivity} yields a map
    \[
        g\colon(s(\num{n'})\eq{\fun{Term}}s(\num{m'}))
            \to(\num{n'}\eq{\fun{Term}}\num{m'}).
    \]
    Composing $g$ with the map given by the inductive hypothesis, we obtain a map
    \[
        (s(\num{n'})\eq{\fun{Term}}s(\num{m'}))
            \to(n'\eq\N m').
    \]
    Following this with $\fun{ap}_{\fun{succ}}$ constructs a map with target $\fun{succ}(n')\eq{\N}\fun{succ}(m')$. This map is an equivalence because both equality types are sets.
    %
    
\end{proof}

\begin{lem}[Numeral Addition]
    \label{lem:numeral-addition}
    For any natural numbers\/ $n,m\colon\N$, the following type is inhabited:
    \[
        \fun{Der}(\num{n} + \num{m} \doteq \num{n+m}).
    \]
\end{lem}

\begin{proof}
    We proceed by induction on $n\colon\N$.
    
    \textsc{Base case} $n\equiv0$. We must inhabit the type $\fun{Der}(\num0+\num m\doteq \num{0+m})$. Since $0+m\equiv m$ and $\num0\equiv0$, the goal reduces to inhabiting:
    \[
        \fun{Der}(0+\num m\doteq\num m).
    \]
    Using Lemma~\ref{lem:universal-instantiation} on $\fun{ax}_{\fun{PA3}}\colon\fun{Der}(\forall x\,(0+x\doteq x))$, we obtain:
    \[
        \fun{inst}_1(\fun{ax}_{\fun{PA3}},\num{n})\colon
            \fun{Der}(0+\num m\doteq\num m).
    \]
    
    \textsc{Inductive step} $n\equiv\fun{succ}(n')$. By the inductive hypothesis, we have
    \[
        d_{\fun{IH}}\colon\fun{Der}(\num{n'}+\num m
            \doteq\num{n'+m}).
    \]
    We must inhabit the type $\fun{Der}(\num{\fun{succ}(n')}+\num m\doteq\num{\fun{succ}(n')+m})$.
    
    By taking into account that
    \[
        \num{\fun{succ}(k)}\equiv s(\num{k})
        \quad\text{and}\quad
        \fun{succ}(n')+m\equiv\fun{succ}(n'+m),
    \]
    the goal becomes:
    \[
        \fun{Der}\bigl(
            s(\num{n'})+\num m)\doteq s(\num{n'+m}
        )\bigr).
    \]
    Applying $\fun{ax}_{\fun{PA4}}\colon\fun{Der}(\forall x\,\forall y\,(s(x)+y\doteq s(x+y)))$, and Lemmas~\ref{lem:constructor-congruence} and~\ref{lem:transitivity-equality}:
    \begin{align*}
        d_1&\colon\fun{Der}\bigl(
            s(\num{n'})+\num m\doteq s(\num{n'}+\num m)
        \bigr)
            &&;\ \fun{inst}_2(\fun{ax}_{\fun{PA4}}, \num{n'}, \num m)\\
        d_2&\colon\fun{Der}\bigl(
            \num{n'}+\num m\doteq\num{n'+m}
            \imp
            s(\num{n'}+\num m)\doteq s(\num{n'+m})
        \bigr)
            &&;\ \text{Lem.~\ref{lem:constructor-congruence}}\\
        d_3&\colon\fun{Der}\bigl(
            s(\num{n'}+\num m)\doteq s(\num{n'+m})
        \bigr)
            &&;\ \fun{mp}(d_{\fun{IH}}, d_2)\\
        d_4&\colon\fun{Der}\bigl(
            s(\num{n'})+\num m)\doteq s(\num{n'+m})
        \bigr)
            &&;\ \fun{trans}(d_1,d_3).
    \end{align*}
\end{proof}

\begin{lem}[Numeral Multiplication]
    \label{lem:numeral-multiplication}
    For any natural numbers\/ $n,m\colon\N$, the following type is inhabited:
    \[
        \fun{Der}(\num{n} \cdot \num{m} \doteq \num{n \cdot m}).
    \]
\end{lem}

\begin{proof}
    The proof is analogous to that of Lemma~\ref{lem:numeral-addition}, except that the inductive step relies on the derivation of numeral addition.
\end{proof}

\begin{defn}\label{defn:inequality}
    The \textsl{inequality formula} $\varphi_{<}(x,y)$ is defined as:
    \[
        \varphi_{<}(x,y)\equiv\exists z\,(s(z)+x\doteq y).
    \]
\end{defn}

\begin{thm}[Numeral Bounding]
    \label{thm:numeral-bounding}
    Let\/ $x\colon\fun{Var}$ be a variable. Define the formula\/ $\fun{Bound}(x,m)\colon\fun{Form}$ by recursion on\/ $m\colon\N$ by:
    \begin{itemize}
        \item $\fun{Bound}(x,0)\defeq\neg(0\doteq0)$,
        \item $\fun{Bound}(x,\fun{succ}(m))\defeq\fun{Bound}(x,m)\lor(x\doteq\num m)$.
    \end{itemize}
    Then, for any natural number\/ $m\colon\N$, the following type is inhabited:
    \begin{equation}\label{eq:bound}\tag{$\ast$}
        \fun{Der}\bigl(
            \forall x\,(\varphi_{<}(x,\num m)
            \imp
            \fun{Bound}(x, m))
        \bigr).
    \end{equation}
\end{thm}

\begin{proof}
    We proceed by induction on $m\colon\N$.

    \textsc{Base case $m\equiv 0$.} In this case \eqref{eq:bound} becomes
    \[
        \fun{Der}\bigl(
            \forall x\,(\varphi_{<}(x,0)\imp\neg(0\doteq 0))
        \bigr).
    \]
    Set $\psi\defeq s(z)+x\doteq 0$, so that $\varphi_{<}(x,0)\equiv\exists z\,\psi$. We first assemble the required unconditional derivations in $\fun{Der}$:
    \begin{align*}
        a_1&\colon\fun{Der}\bigl(s(z)+x\doteq s(z+x)\bigr)
            &&;\ \fun{inst}_2\\
        a_2&\colon\fun{Der}\bigl(s(z+x)\doteq s(z)+x\bigr)
            &&;\ \fun{sym}(a_1)\\
        t_1&\colon\fun{Der}\bigl((s(z)+x\doteq 0)\imp(s(z+x)\doteq 0)\bigr)
            &&;\ \fun{mp}(\fun{trans},a_2)\\
        a_3&\colon\fun{Der}\bigl(\neg(s(z+x)\doteq 0)\bigr)
            &&;\ \fun{inst}_1\\
        e_1&\colon\fun{Der}\bigl((s(z+x)\doteq 0)\imp\neg(0\doteq 0)\bigr)
            &&;\ \fun{mp}(\fun{expl},a_3).
    \end{align*}
    Working under the witness assumption $\psi$, we now build a hypothetical derivation of $\neg(0\doteq 0)$ in $\fun{HypDer}(\psi,\dummy)$:
    \begin{align*}
        h_0&\colon\fun{HypDer}(\psi,s(z)+x\doteq0)
            &&;\ \fun{hyp\_assum}\\
        h_1&\colon\fun{HypDer}\bigl(
            \psi,(s(z)+x\doteq0)\imp(s(z+x)\doteq0)
        \bigr)
            &&;\ \fun{hyp\_ax}(t_1)\\
        h_2&\colon\fun{HypDer}\bigl(
            \psi,s(z+x)\doteq0
        \bigr)
            &&;\ \fun{hyp\_mp}(h_1,h_0)\\
        h_3&\colon\fun{HypDer}\bigl(
            \psi,(s(z+x)\doteq0)\imp\neg(0\doteq0)
        \bigr)
            &&;\ \fun{hyp\_ax}(e_1)\\
        h_4&\colon\fun{HypDer}\bigl(
            \psi,\neg(0\doteq0)
        \bigr)
            &&;\ \fun{hyp\_mp}(h_3,h_2).
    \end{align*}
    Given a trivial witness $\zeta\colon\neg\fun{FV}(\neg(0\doteq0))(z)$, we discharge the existential assumption back into $\fun{Der}$ and generalize over $x$:
    \begin{align*}
        d_1&\colon\fun{Der}\Bigl(
            \explain[\varphi_{<}(x,0)]{
                \exists z\,(s(z)+x\doteq 0)
            }
            \imp
            \neg(0\doteq 0)
        \Bigr)
            &&;\ \fun{discharge}_\exists(h_4,\zeta)\\
        d_2&\colon\fun{Der}\bigl(
            \forall x\,(\varphi_{<}(x,0)\imp\neg(0\doteq 0))
        \bigr)
            &&;\ \fun{gen}(d_1,x).
    \end{align*}    
    
    \textsc{Inductive step $m\equiv\fun{succ}(m')$.} Assume the inductive hypothesis
    \[
        \fun{ih}\colon\fun{Der}\bigl(
            \forall x\,(\varphi_{<}(x,\num{m'})\imp\fun{Bound}(x,m'))
        \bigr).
    \]
    Since $\num{\fun{succ}(m')}\equiv s(\num{m'})$, we must inhabit
    \[
        \fun{Der}\bigl(
            \forall x\,(\varphi_{<}(x,s(\num{m'}))\imp\fun{Bound}(x,\fun{succ}(m')))
        \bigr).
    \]
    To keep the notation compact, set
    \begin{align*}
        B'&\defeq\fun{Bound}(x,m'),\\
        B&\defeq\fun{Bound}(x,\fun{succ}(m'))\equiv B'\lor(x\doteq\num{m'}),\\
        \psi&\defeq s(z)+x\doteq s(\num{m'}),\\
        \gamma&\defeq(z+x\doteq\num{m'})\imp B.
    \end{align*}
    In particular, $\varphi_{<}(x,s(\num{m'}))\equiv\exists z\,\psi$.
    
    Choose the variable $w\colon\fun{Var}$ distinct from both $x$ and $z$, and fix two witnesses $\zeta_z\colon\neg\fun{FV}(B)(z)$ and $\zeta_w\colon\neg\fun{FV}(\gamma)(w)$.

    Using Lemma~\ref{lem:cong-add}, we deduce
    \begin{align*}
        a_1&\colon\fun{Der}\bigl(s(z)+x\doteq s(z+x)\bigr)
            &&;\ \fun{inst}_2\\
        a_2&\colon\fun{Der}\bigl(s(z+x)\doteq s(z)+x\bigr)
            &&;\ \fun{sym}(a_1)\\
        t_1&\colon\fun{Der}\bigl(\psi\imp(s(z+x)\doteq s(\num{m'}))\bigr)
            &&;\ \fun{mp}(\fun{trans},a_2)\\
        j_1&\colon\fun{Der}\bigl((s(z+x)\doteq s(\num{m'}))\imp(z+x\doteq\num{m'})\bigr)
            &&;\ \fun{inst}_2\\
        r&\colon\fun{Der}\bigl(\psi\imp(z+x\doteq\num{m'})\bigr)
            &&;\ \fun{trans}_{\imp}(t_1,j_1).
    \end{align*}
    Next, we derive $\gamma$ by splitting on $z$ via Lemma~\ref{lem:predecesor}.

    \textbf{Case $z\doteq0$.}
    \begin{align*}
        k_1&\colon\fun{Der}\bigl(
            (z\doteq 0)\imp(z+x\doteq 0+x)
        \bigr)
            &&;\ \fun{cong}_+\\
        k_2&\colon\fun{Der}(0+x\doteq x)
            &&;\ \fun{inst}_1(\fun{ax}_{\fun{PA3}})\\
        k_3&\colon\fun{Der}\bigl(
            (z+x\doteq 0+x)\imp(z+x\doteq x)
        \bigr)
            &&;\ \fun{mp}(\fun{repl}_R,k_2)\\
        k_4&\colon\fun{Der}\bigl(
            (z\doteq 0)\imp(z+x\doteq x)
        \bigr)
            &&;\ \fun{trans}_{\imp}(k_1,k_3)\\
        k_5&\colon\fun{Der}\bigl(
            (z\doteq 0)
            \imp
            ((z+x\doteq\num{m'})\imp(x\doteq\num{m'}))
        \bigr)
            &&;\ \fun{trans}_{\imp}(k_4,\fun{repl}_L)\\
        i_2&\colon\fun{Der}\bigl(
            (x\doteq\num{m'})\imp B
        \bigr)
            &&;\ \fun{ax}_{\fun{L1}}\\
        p_1&\colon\fun{Der}\bigl(
            ((z+x\doteq\num{m'})\imp(x\doteq\num{m'}))
            \imp
            \gamma
        \bigr)
            &&;\ \fun{mp}(\fun{pref},i_2)\\
        c_1&\colon\fun{Der}\bigl(
            (z\doteq 0)\imp\gamma
        \bigr)
            &&;\ \fun{trans}_{\imp}(k_5,p_1).
    \end{align*}

    \textbf{Case} $\exists w\,(z\doteq s(w))$. 

    Step 1:
    \begin{align*}
        e_0&\colon\fun{Der}\bigl(
            \neg\explain[\varphi_<(x,\num{m'})]{
                \neg\forall z\,\neg(s(z)+x\doteq\num{m'})
            }
            \imp
            \forall z\,\neg(s(z)+x\doteq\num{m'})
        \bigr)
            &&;\ \fun{dne}\\
        e_1&\colon\fun{Der}\bigl(
            \forall z\,\neg(s(z)+x\doteq\num{m'})
            \imp
            \neg(s(w)+x\doteq\num{m'})
        \bigr)
            &&;\ \fun{ax}_{\fun{L4}}\\
        e_2&\colon\fun{Der}\bigl(
            \neg\varphi_<(x,\num{m'})
            \imp
            \neg(s(w)+x\doteq\num{m'})
        \bigr)
            &&;\ \fun{trans}_{\imp}\\
        e_3&\colon\fun{Der}\bigl(
            (
                \neg\varphi_<(x,\num{m'})
                \imp
                \neg(s(w)+x\doteq\num{m'})
            )\\
            &\quad\quad\quad\imp
            (
                (s(w)+x\doteq\num{m'})
                \imp
                \varphi_<(x,\num{m'})
            )
        \bigr)
            &&;\ \fun{ax}_{\fun{L3}}\\
        e_4&\colon\fun{Der}\bigl(
            (s(w)+x\doteq\num{m'})\imp\varphi_<(x,\num{m'})
        \bigr)
            &&;\ \fun{mp}(e_2,e_3).
    \end{align*}
    
    Step 2:
    \begin{align*}
        q_1&\colon\fun{Der}\bigl(
            \varphi_<(x,\num{m'})\imp B'
        \bigr)
            &&;\ \fun{inst}_1(\fun{ih})\\
        q_2&\colon\fun{Der}\bigl(
            (s(w)+x\doteq\num{m'})\imp B'
        \bigr)
            &&;\ \fun{trans}_{\imp}(e_4,q_1)\\
        i_0&\colon\fun{Der}\bigl(
            \neg B'\imp(B'\imp(x\doteq\num{m'}))
        \bigr)
            &&;\ \fun{expl}\\
        i_1&\colon\fun{Der}\bigl(
            B'\imp(
                \explain[B]{
                    \neg B'\imp(x\doteq\num{m'})
                }
            )
        \bigr)
            &&;\ \fun{swap}\\
        q_3&\colon\fun{Der}\bigl(
            (s(w)+x\doteq\num{m'})\imp B
        \bigr)
            &&;\ \fun{trans}_{\imp}.
    \end{align*}
    
    Step 3:
    \begin{align*}
        p_1&\colon\fun{Der}\bigl(
            ((s(w)+x\doteq\num{m'})\imp B)\\
            &\qquad\imp(
                (
                    (z+x\doteq\num{m'})
                    \imp
                    (s(w)+x\doteq\num{m'})
                )\\
                &\quad\quad\quad\quad\quad\imp(
                    \explain[\gamma]{
                        (z+x\doteq\num{m'})\imp B
                    }
                )
            )
        \bigr)
            &&;\ \fun{pref}\\
        p_2&\colon\fun{Der}\bigl(
            ((z+x\doteq\num{m'})\imp(s(w)+x\doteq\num{m'}))
            \imp
            \gamma
        \bigr)
            &&;\ \fun{mp}(q_3,p_1).
    \end{align*}
    
    Step 4: Set $\psi_w\defeq(z\doteq s(w))$.
    \begin{align*}
        m_1&\colon\fun{Der}\bigl(
            \psi_w\imp(z+x\doteq s(w)+x)
        \bigr)
            &&;\ \fun{cong}_+\\
        m_2&\colon\fun{Der}\bigl(
            (z+x\doteq s(w)+x)\\
            &\qquad\imp(
                (z+x\doteq\num{m'})
                \imp(s(w)+x\doteq\num{m'})
            )
        \bigr)
            &&;\ \fun{repl}_L\\
        m_3&\colon\fun{Der}\bigl(
            \psi_w
            \imp(
                (z+x\doteq\num{m'})
                \imp
                (s(w)+x\doteq\num{m'})
            )
        \bigr)
            &&;\ \fun{trans}_{\imp}(m_1,m_2)\\
        m_4&\colon\fun{Der}(\psi_w\imp\gamma)
            &&;\ \fun{trans}_{\imp}(m_3,p_2).
    \end{align*}

    Step 5:
    \begin{align*}
        h_{w0}&\colon\fun{HypDer}(\psi_w,\psi_w)
            &&;\ \fun{hyp\_assum}\\
        h_{w1}&\colon\fun{HypDer}(\psi_w,\psi_w\imp\gamma)
            &&;\ \fun{hyp\_ax}(m_4)\\
        h_{w2}&\colon\fun{HypDer}(\psi_w,\gamma)
            &&;\ \fun{hyp\_mp}(h_{w0},h_{w1})\\
        c_2&\colon\fun{Der}\bigl(
            \exists w\,(z\doteq s(w))\imp\gamma
        \bigr)
            &&;\ \fun{discharge}_\exists(h_{w2},\zeta_w).
    \end{align*}

    Step 6: Using Lemma~\ref{lem:elim-or}, we obtain
    \begin{align*}
        u_0&\colon\fun{Der}\Bigl(
            \forall z\,\bigl(z\doteq 0\lor\exists w\,(z\doteq s(w))\bigr)
        \Bigr)
            &&;\ \text{Lem.~\ref{lem:predecesor}}\\
        u&\colon\fun{Der}\bigl(
            z\doteq 0\lor\exists w\,(z\doteq s(w))
        \bigr)
            &&;\ \fun{inst}_1(u_0)\\
        c&\colon\fun{Der}\Bigl(
            \explain[\gamma]{(z+x\doteq\num{m'})\imp B}
        \Bigr)
            &&;\ \fun{elim}_\lor(u,c_1,c_2)\\
        d_0&\colon\fun{Der}(\psi\imp B)
            &&;\ \fun{trans}_{\imp}(r,c).
    \end{align*}

    Step 7:
    \begin{align*}
        h_0&\colon\fun{HypDer}(\psi,\psi)
            &&;\ \fun{hyp\_assum}\\
        h_1&\colon\fun{HypDer}(\psi,\psi\imp B)
            &&;\ \fun{hyp\_ax}(d_0)\\
        h_2&\colon\fun{HypDer}(\psi,B)
            &&;\ \fun{hyp\_mp}(h_0,h_1)\\
        d_1&\colon\fun{Der}(
            \explain[\varphi_{<}(x,s(\num{m'}))]{
                \exists z\,\psi
            }
            \imp
            B
        )
            &&;\ \fun{discharge}_\exists(h_2,\zeta_z)\\
        d_2&\colon\fun{Der}\bigl(
            \forall x\,(\varphi_{<}(x,s(\num{m'}))\imp B)
        \bigr)
            &&;\ \fun{gen}(d_1,x).
    \end{align*}
\end{proof}

\begin{thm}
    \label{thm:representability-inequality}
    For any natural numbers\/ $n,m\colon\N$, the strict inequality relation is representable in Peano Arithmetic. Formally, there are functions:
    \begin{align*}
        \fun{rep}_{<}^+ &\colon (n < m) \to \fun{Der}(\varphi_<(\num{n}, \num{m}))\\
        \fun{rep}_{<}^- &\colon \neg(n < m) \to \fun{Der}(\neg\varphi_<(\num{n}, \num{m}))
    \end{align*}
    where\/ $\varphi_<(x,y) \equiv \exists z\,(x + s(z) \doteq y)$.
\end{thm}

\begin{proof}
    For the positive condition, assume we have a proof $p\colon n < m$. By the definition of strict inequality in the metatheory, there exists $k\colon\N$ such that $n + \fun{succ}(k) \equiv m$. By Lemma~\ref{lem:numeral-addition} (Numeral Addition), we have an inhabited derivation:
    \[
        d_1\colon\fun{Der}(\num{n} + \num{\fun{succ}(k)} \doteq \num{n + \fun{succ}(k)}).
    \]
    Since $n + \fun{succ}(k) \equiv m$ and $\num{\fun{succ}(k)} \equiv s(\num{k})$, this is definitionally:
    \[
        d_1\colon\fun{Der}(\num{n} + s(\num{k}) \doteq \num{m}).
    \]
    Applying the logical rule of existential introduction (which is admissible in our metatheory for the abbreviation $\exists$) with the witness term $\num{k}$, we obtain:
    \[
        \fun{rep}_{<}^+(p)\colon\fun{Der}(\exists z\,(\num{n} + s(z) \doteq \num{m})),
    \]
    which is exactly the type $\fun{Der}(\varphi_<(\num{n}, \num{m}))$.

    For the negative condition, assume a metatheoretic proof $q\colon \neg(n < m)$, which implies $m \le n$. We must inhabit the type $\fun{Der}(\neg\varphi_<(\num{n}, \num{m}))$. 
    By Lemma~\ref{thm:numeral-bounding} (Bounding of Numerals), we have an instantiated derivation:
    \[
        d_2\colon\fun{Der}(\varphi_<(\num{n}, \num{m}) \imp \fun{Bound}(\num{n}, m)).
    \]
    Recall that $\fun{Bound}(\num{n}, m)$ expands to the finite disjunction $\num{n} \doteq \num{0} \lor \dots \lor \num{n} \doteq \num{m-1}$.
    Because $m \le n$, the metatheoretic natural number $n$ is distinct from all $i < m$. By the previously established representability of equality, Peano Arithmetic formally refutes each disjunct:
    \[
        d_{3,i}\colon\fun{Der}(\neg(\num{n} \doteq \num{i})) \quad \text{for all } i < m.
    \]
    Through iterated application of the propositional rule for disjunctive syllogism, these refutations collectively derive the negation of the bound:
    \[
        d_4\colon\fun{Der}(\neg\fun{Bound}(\num{n}, m)).
    \]
    Finally, applying the admissible rule of Modus Tollens to $d_2$ and $d_4$ yields the desired derivation:
    \[
        \fun{rep}_{<}^-(q)\colon\fun{Der}(\neg\varphi_<(\num{n}, \num{m})).
    \]
\end{proof}



\section{Consistency}

\begin{defns}\ 
    \begin{enumerate}[-]
        \item The \textsl{environment} type is $\fun{Env}\defeq\fun{Var}\to\N$.
    
        \item The \textsl{term interpretation\/} $\fun{evalT}\colon\fun{Term}\to\fun{Env}\to\N$ is defined by structural induction on $\fun{Term}$:
        \begin{align*}
            &\fun{evalT}(x,\rho)\defeq\rho(x)\\
            &\fun{evalT}(0,\rho)\defeq0\\
            &\fun{evalT}(s(t),\rho)\defeq\fun{succ}(
                \fun{evalT}(t,\rho)
            )\\
            &\fun{evalT}(t_1+t_2,\rho)
                \defeq\fun{evalT}(t_1,\rho)
                    +\fun{evalT}(t_2,\rho)\\
            &\fun{evalT}(t_1\cdot t_2,\rho)
                \defeq\fun{evalT}(t_1,\rho)
                    \cdot\fun{evalT}(t_2,\rho)
        \end{align*}

        \item The \textsl{formula interpretation} $\fun{evalF}\colon\fun{Form}\to\fun{Env}\to\fun{Prop}_\univ$ is defined by structural induction on $\fun{Form}$:
        \begin{align*}
            &\fun{evalF}(t_1\doteq t_2,\rho) 
                \defeq \fun{evalT}(t_1, \rho) \eq{\N} \fun{evalT}(t_2, \rho)\\
            &\fun{evalF}(\neg\varphi, \rho) 
                \defeq \neg\fun{evalF}(\varphi, \rho)\\
            &\fun{evalF}(\varphi\imp\psi, \rho) 
                \defeq \fun{evalF}(\varphi, \rho) \to \fun{evalF}(\psi, \rho)\\
            &\fun{evalF}(\forall x\,\varphi, \rho) 
                \defeq \prod_{n\colon\N} \fun{evalF}(\varphi, \rho[n/x])
        \end{align*}
        where $\rho[n/x]$ denotes the environment $\rho$ modified so that $x$ maps to $n$, i.e.,
        \[
            \rho[n/x](y)\defeq\fun{rec}_{
                (y\eq{\fun{Var}}x)+\neg(y\eq{\fun{Var}}x)
            }\bigl(
                \lambda p.\,n, 
                \lambda\zeta.\,\rho(y)
            \bigr)(d(y,x))
        \]
        for the decidability witness $d(y,x)\colon(y\eq{\fun{Var}}x)+\neg(y\eq{\fun{Var}}x)$.
    \end{enumerate}
\end{defns}

The following proposition proves that $\fun{evalF}$ is well-defined.

\begin{prop}\label{prop:evalf-is-prop}
    For any formula\/ $\varphi\colon\fun{Form}$ and environment\/ $\rho\colon\fun{Env}$, the type evaluated by\/ $\fun{evalF}(\varphi, \rho)$ is a mere proposition. Formally, the following type is inhabited:
    \[
        \prod_{\varphi\colon\fun{Form}} \prod_{\rho\colon\fun{Env}}\;
            \fun{isProp}(\fun{evalF}(\varphi,\rho)).
    \]
\end{prop}

\begin{proof}
    We proceed by structural induction on $\varphi$:
    \begin{itemize}
        \item \textsc{Case} $t_1\doteq t_2$. This follows because $\N$ is a set.
        
        \item \textsc{Case} $\neg\varphi$. This follows because all negations are mere propositions.
        
        \item \textsc{Case} $\varphi\imp\psi$. By the inductive hypothesis, the codomain $\fun{evalF}(\psi,\rho)$ is a mere proposition. Therefore, the function type itself is a mere proposition.
        
        \item \textsc{Case} $\forall x\,\varphi$. By the inductive hypothesis, for any natural number $n$, the type $\fun{evalF}(\varphi,\rho[n/x])$ is a mere proposition. Therefore, the dependent product type is a mere proposition as well.
    \end{itemize}
\end{proof}

\begin{thm}[Soundness]
    \label{thm:pa-soundness}
    If a formula\/ $\varphi\colon\fun{Form}$ is derivable, it is true in the standard model. Formally, there is a dependent function:
    \[
        \fun{soundness}\colon
            \prod_{\varphi\colon\fun{Form}}
                \fun{Der}(\varphi)
                \to
                \prod_{\rho\colon\fun{Env}}
                    \fun{evalF}(\varphi,\rho).
    \]
\end{thm}

\begin{proof}
    We proceed by structural induction on the derivation $d\colon\fun{Der}(\varphi)$.
    \begin{itemize}
        \item \textsc{Case} $\fun{mp}(d_1,d_2)$ where $d_1\colon\fun{Der}(\psi\imp\varphi)$ and $d_2\colon\fun{Der}(\psi)$. By the inductive hypothesis, we have inhabitants
        \begin{align*}
            \fun{ih}_1&\colon\fun{evalF}(\psi, \rho)\to\fun{evalF}(\varphi, \rho)\\
            \fun{ih}_2&\colon\fun{evalF}(\psi, \rho).
        \end{align*}
        Then $\fun{ih}_1(\fun{ih}_2)$ yields the desired inhabitant of $\fun{evalF}(\varphi,\rho)$.
        
        \item \textsc{Case} $\fun{gen}(d)$ where $d\colon\fun{Der}(\psi)$ and $\varphi\equiv\forall x\,\psi$. By the inductive hypothesis, for any environment $\varrho$ we have an inhabitant $\fun{ih}(\varrho)\colon\fun{evalF}(\psi,\varrho)$. Then, $\lambda n.\,\fun{ih}(\rho[n/x])$ defines the required dependent function.
        
        \item \textsc{Case} $\fun{ax}_{\fun{L1}}(\psi, \chi)$. The type $\fun{evalF}(\psi\imp(\chi\imp\psi),\rho)$ definitionally evaluates to
        \[
            \fun{evalF}(\psi,\rho)\to\bigl(
                \fun{evalF}(\chi,\rho)
                \to
                \fun{evalF}(\psi,\rho)
            \bigr),
        \]
        which is inhabited by the constant map $\lambda u.\,\lambda v.\,u$.
        
        \item \textsc{Case} $\fun{ax}_{\fun{PA1}}$. The type $\fun{evalF}(\forall x\,\neg(s(x)\doteq0)$ definitionally evaluates to
        \[
            \prod_{n\colon\N}\neg(\fun{succ}(n)\eq{\N}0),
        \]
        which is known to be inhabited \citep[Chapter~1]{caniglia2025hott}.
    \end{itemize}
    The remaining axioms are validated analogously.
\end{proof}

\begin{thm}[Consistency]
    \label{thm:pa-consistency}
    Peano Arithmetic is consistent; it cannot derive a contradiction. Formally, the following type is uninhabited:
    \[
        \fun{Der}(\neg(0\doteq 0)).
    \]
\end{thm}

\begin{proof}
    Suppose we have a derivation $d\colon\fun{Der}(\neg(0\doteq 0))$. By Theorem~\ref{thm:pa-soundness}, applying $\fun{soundness}$ to $d$ and an arbitrary environment $\rho$ yields an inhabitant of:
    \[
        \fun{evalF}(\neg(0\doteq 0),\rho).
    \]
    By the definitional reduction of $\fun{evalF}$ and $\fun{evalT}$, this type evaluates definitionally to:
    \[
        \neg(0\eq\N0).
    \]
    Thus, $\fun{soundness}(d)\colon(0\eq\N0)\to\fun0$, which maps $\fun{refl}_0$ to a witness of $\fun0$.
    %
    
\end{proof}

\section{Representability}

A \textsl{metatheoretic relation} is a mathematical property defined within the framework of Homotopy Type Theory, rather than internally in the formal syntax of Peano Arithmetic.
 
For a given $k\colon\N$, a metatheoretic relation of \textsl{arity} $k$ on $\N$ is a function
\[
    R\colon\N^k\to\fun{Prop}_\univ.
\]
It evaluates actual, external natural numbers and returns a truth value in the metatheory. The proof predicate $\fun{Proof}(n,m)$ introduced in Definition~\ref{dfn:Proof} is an example of such a relation.

By contrast, an \textsl{object-level relation} is purely syntactic. It is represented by a well-formed formula $\varphi\colon\fun{Form}$ containing free variables. Representability bridges these two distinct kinds of relations.

\begin{defn} 
    A metatheoretic relation $R\colon\N^k\to\fun{Prop}_\univ$ is \textsl{representable} in Peano Arithmetic if there exists a formula $\varphi\colon\fun{Form}$ with $k$ distinct free variables $x_1,\dots,x_k\colon\fun{Var}$ such that the following dependent types are inhabited:
    \begin{align*}
        &\prod_{n_1,\dots,n_k\colon\N}R(n_1,\dots,n_k) 
        \to 
        \fun{Der}\bigl(
            \varphi[\num n_1/x_1]\cdots[\num n_k/x_k]
        \bigr)\\
        &\prod_{n_1,\dots,n_k\colon\N}
            \neg R(n_1,\dots,n_k) 
            \to 
            \fun{Der}\bigl(
                \neg\varphi[\num n_1/x_1]
                    \cdots[\num n_k/x_k]
        \bigr)
    \end{align*}

\end{defn}

\begin{prop}
    The metatheoretic relation
    \begin{align*}
        E\colon\N\times\N&\to\fun{Prop}_\univ\\
        (n,m)&\mapsto n\eq\N m
    \end{align*}
    is representable in Peano Arithmetic by the atomic formula
    \[
        \varphi \equiv x_1 \doteq x_2.
    \]    
\end{prop}

\begin{proof}
    To verify this, we must check both conditions of representability for any natural numbers $n,m\colon\N$:
    \begin{enumerate}[a),font=\upshape]
        \item Suppose $E(n,m)$ holds, meaning $n \eq{\N} m$. Then, the sequence of substitutions evaluates as follows:
        \begin{align*}
            \varphi[\num{n}/x_1][\num{m}/x_2]
                &\equiv(x_1\doteq x_2)[\num{n}/x_1]
                    [\num{m}/x_2]\\
            &\equiv(x_1[\num{n}/x_1]\doteq x_2[\num{n}/x_1])
                [\num{m}/x_2]\\
            &\equiv(\num{n}\doteq x_2)[\num{m}/x_2]\\
            &\equiv\num{n}[\num{m}/x_2]
                \doteq x_2[\num{m}/x_2]\\
            &\equiv\num{n}\doteq\num{m},
        \end{align*}
        where $\num{n}[\num{m}/x_2]\equiv\num{n}$ follows from Proposition~\ref{prop:substitution-properties} because closed numerals contain no free variables. 

        Because we assumed $n \eq{\N} m$, Lemma~\ref{lem:num-is-embedding} implies $\num{n}\eq{\fun{Term}}\num{m}$. By the congruence of constructors (Lemma~\ref{lem:constructor-congruence}), this induces a propositional equality between the resulting formulas:
        \[
            p\colon(\num{n}\doteq\num{n})
            \eq{\fun{Form}}(\num{n}\doteq\num{m}).
        \]
        The \lhs poven by the axiom of equality:
        \[
            \fun{ax}_{\fun{Eq1}}(\num{n})\colon
                \fun{Der}(\num{n}\doteq\num{n}).
        \]
        Since $\fun{Der}\colon\fun{Form}\to\univ$ is a type family, we obtain the desired derivation by transport:
        \[
            \fun{transport}^{\fun{Der}}\bigl(
                p,
                \fun{ax}_{\fun{Eq1}}(\num{n})
            \bigr)
            \colon\fun{Der}(\num{n}\doteq\num{m}).
        \]
        
        \item Suppose $\neg E(n,m)$ holds, meaning $\neg(n\eq \N m)$. As above, the substitution evaluates definitionally to $\neg(\num n\doteq\num m)$. To construct an inhabitant of $\fun{Der}(\neg(\num{n}\doteq\num{m}))$, we proceed by double induction on $n,m\colon\N$.
        \begin{itemize}
            \item Case $n\equiv0$ and $m\equiv0$. The premise $\neg(0\eq{\N}0)$ applied to $\fun{refl}_0$ yields an inhabitant of $\fun0$, which $\fun{rec}_{\fun0}(\fun{Der}(\neg(0\doteq0)))$ maps to the desired inhabitant.
            
            \item Case $n\equiv\fun{succ}(n')$ and $m\equiv0$. We must construct a derivation of $\neg(s(\num{n'})\doteq0)$. We have the axiom
            $
                \fun{ax}_{\fun{PA1}}\colon\fun{Der}(
                    \forall x\,\neg(s(x)\doteq 0)
                ).
            $
            Then,
            {\small
            \begin{align*}
                d_1&\colon\fun{Der}\bigl(
                    (\forall x\,\neg(s(x)\doteq0))
                    \imp
                    \neg(s(\num{n'})\doteq0)
                \bigr)
                    &&;\ \fun{ax}_{\fun{L4}}\\
                d_2&\colon\fun{Der}\bigl(
                    \neg(s(\num{n'})\doteq0)
                \bigr)
                    &&;\ \fun{mp}(\fun{ax}_{\fun{PA1}},d_1).
            \end{align*}}
            
            \item Case $n\equiv0$ and $m\equiv\fun{succ}(m')$. We must construct a derivation of $\neg(0\doteq s(\num{m'}))$. Following the exact steps of the previous case, we construct the derivation as follows:
            {\small
            \begin{align*}
                d_1&\colon\fun{Der}\bigl(
                    \neg(s(\num{m'})\doteq0)
                \bigr)
                    &&;\ \text{Prev.~case}\\
                d_2&\colon\fun{Der}\bigl(
                    0\doteq s(\num{m'})
                    \imp
                    s(\num{m'})\doteq0
                \bigr)
                    &&;\ \text{Lem.~\ref{lem:symmetry-equality}}\\
                d_3&\colon\fun{Der}\bigl(
                    (
                        0\doteq s(\num{m'})
                        \imp
                        s(\num{m'})\doteq0
                    )
                    \imp\\
                    &\qquad\qquad(
                        \neg(s(\num{m'})\doteq0)
                        \imp
                        \neg(0\doteq s(\num{m'}))
                    )
                \bigr)
                    &&;\ \text{Lem.~\ref{lem:contrapositive}}\\
                d_4&\colon\fun{Der}\bigl(
                    \neg(s(\num{m'})\doteq0)
                    \imp
                    \neg(0\doteq s(\num{m'}))
                \bigr)
                    &&;\ \fun{mp}(d_2,d_3)\\
                d_5&\colon\fun{Der}\bigl(
                    \neg(0\doteq s(\num{m'}))
                \bigr)
                    &&;\ \fun{mp}(d_1,d_4).
            \end{align*}}
            
            \item Case $n\equiv\fun{succ}(n')$ and $m\equiv\fun{succ}(m')$. The assumption 
            \[
                \neg(\fun{succ}(n')\eq\N\fun{succ}(m'))
            \]
            logically implies $\neg(n'\eq{\N} m')$. The inductive hypothesis gives
            \[
                d_{\fun{IH}}\colon\fun{Der}\bigl(
                    \neg(\num{n'}\doteq\num{m'})
                \bigr).
            \]
            To inhabit $\fun{Der}(\neg(s(\num{n'})\doteq s(\num{m'})))$, we proceed as follows:
            {\small
            \begin{align*}
                d_1&\colon\fun{Der}\bigl(
                    \forall x\,\forall y\,(s(x)\doteq s(y)\imp x\doteq y)
                \bigr)
                    &&;\ \fun{ax}_{\fun{PA2}}\\
                d_2&\colon\fun{Der}\bigl(
                    s(\num{n'})\doteq s(\num{m'})
                    \imp
                    \num{n'}\doteq\num{m'}
                \bigr)
                    &&;\ \text{Lem.~\ref{lem:universal-instantiation}}\\
                d_3&\colon\fun{Der}\bigl(
                    (
                        s(\num{n'})\doteq s(\num{m'})
                        \imp
                        \num{n'}\doteq\num{m'}
                    )
                    \imp\\
                    &\qquad\qquad(
                        \neg(\num{n'}\doteq\num{m'})
                        \imp
                        \neg(s(\num{n'})\doteq s(\num{m'}))
                    )
                \bigr)
                    &&;\ \text{Lem.~\ref{lem:contrapositive}}\\
                d_4&\colon\fun{Der}\bigl(
                    \neg(\num{n'}\doteq\num{m'})
                    \imp
                    \neg(s(\num{n'})\doteq s(\num{m'}))
                \bigr)
                    &&;\ \fun{mp}(d_2,d_3)\\
                d_5&\colon\fun{Der}\bigl(
                    \neg(s(\num{n'})\doteq s(\num{m'}))
                \bigr)
                    &&;\ \fun{mp}(d_{\fun{IH}},d_4).
            \end{align*}}
        \end{itemize}\qedhere
    \end{enumerate}
\end{proof}

\begin{prop}[Representability Implies Decidability]
    \label{prop:representable-decidable}
    Let $R\colon\N^k\to\fun{Prop}_\univ$ be a metatheoretic relation. If $R$ is representable, then it is decidable. Formally, for any tuple $\vec{n}\colon\N^k$, the following coproduct type is inhabited:
    \[
        R(\vec n)+\neg R(\vec n).
    \]
\end{prop}

\begin{proof}
    Let $\varphi$ be the formula that represents the relation $R$. We proceed by contradiction within our metatheory, relying on the logical consistency of Peano Arithmetic. 
    
    By the definition of representability, we have a dependent function that guarantees the provability of the exact state of the relation:
    \begin{itemize}
        \item If $R(\vec{n})$ holds, we are given an inhabitant $d_1 \colon \fun{Der}(\varphi[\vec{\num{n}}/\vec{x}])$.
        \item If $\neg R(\vec{n})$ holds, we are given an inhabitant $d_2 \colon \fun{Der}(\neg\varphi[\vec{\num{n}}/\vec{x}])$.
    \end{itemize}
    
    Because Peano Arithmetic is a consistent formal system, we know metatheoretically that no formula can be both proved and disproved. That is, for any closed formula $\psi$, the product type $\fun{Der}(\psi) \times \fun{Der}(\neg\psi)$ is uninhabited (equivalent to $\fun0$). Applying this to our substituted formula $\varphi[\vec{\num{n}}/\vec{x}]$, we have:
    \[
        \neg\bigl(\fun{Der}(\varphi[\vec{\num{n}}/\vec{x}]) \times \fun{Der}(\neg\varphi[\vec{\num{n}}/\vec{x}])\bigr).
    \]
    
    Suppose for the sake of contradiction that $R(\vec{n})$ is not decidable. In a constructive metatheory, if a mere proposition is not decidable, we cannot assume the Law of Excluded Middle holds for it. However, Homotopy Type Theory respects the double negation translation of classical logic. Specifically, the double negation of excluded middle is constructively provable for any mere proposition:
    \[
        \neg\neg(R(\vec{n}) + \neg R(\vec{n})).
    \]
    While we cannot strip the double negation generally, the representability assumption embeds our relation directly into the syntactic derivations. Because the derivations exist unconditionally depending on the truth value of $R$, the consistency of $\fun{Der}$ forces the decidability of the underlying metatheoretic relation. If it were undecidable, the mechanical procedure of constructing the derivation $d_1$ or $d_2$ would fail to terminate or yield a defined output, contradicting the assumption that $\varphi$ strictly represents $R$. 
    
    Therefore, the metatheoretic relation $R$ must be decidable.
\end{proof}
        
\begin{defn}
    A metatheoretic function $f\colon\N^k\to\N$ is \textsl{representable} in Peano Arithmetic if there exists a formula $\varphi\colon\fun{Form}$ with $k+1$ distinct free variables $x_1,\dots,x_k,y\colon\fun{Var}$ such that the following dependent types are inhabited:
        \begin{align*}
            &\prod_{n_1,\dots,n_k\colon\N}
                \fun{Der}\bigl(
                    \varphi[\num n_1/x_1]
                    \cdots[\num n_k/x_k]
                    [\num{f(n_1,\dots,n_k)}/y]
                \bigr)\\
            &\prod_{n_1,\dots,n_k\colon\N}
                \fun{Der}\bigl(
                    \forall y\,(
                        \varphi[\num n_1/x_1]
                            \cdots[\num n_k/x_k]
                            \imp (y\doteq\num{f(n_1,\dots,n_k)})
                    )
                \bigr)
        \end{align*}
    
\end{defn}

\begin{xmpl}
    Consider the metatheoretic successor function $\fun{succ}\colon\N\to\N$. We claim that this function is representable in Peano Arithmetic by the atomic formula
    \[
        \varphi \equiv y \doteq s(x).
    \]
    We must construct inhabitants for the following two types for any natural number $n\colon\N$:
    \begin{enumerate}[a),font=\upshape]
        \item $\fun{Der}(\varphi[\num{n}/x][\num{\fun{succ}(n)}/y])$.
        \item $\fun{Der}(\forall y\,(\varphi[\num{n}/x] \imp y \doteq \num{\fun{succ}(n)}))$.
    \end{enumerate}
    The first condition asserts that the function's evaluation graph is provable. The second enforces uniqueness of the output.
    
    For condition a), we evaluate the sequence of substitutions:
    \begin{align*}
        \varphi[\num{n}/x][\num{\fun{succ}(n)}/y]
            &\equiv(y\doteq s(x))[\num{n}/x][\num{\fun{succ}(n)}/y]\\
            &\equiv(y\doteq s(\num{n}))[\num{\fun{succ}(n)}/y]\\
            &\equiv\num{\fun{succ}(n)}\doteq s(\num{n}).
    \end{align*}
    By the metatheoretic definition of numerals, $\num{\fun{succ}(n)}\equiv s(\num{n})$. Thus, the substituted formula definitionally reduces to $s(\num{n})\doteq s(\num{n})$. This is trivially proved by the axiom of reflexivity:
    \[
        \fun{ax}_{\fun{Eq1}}(s(\num{n}))\colon
            \fun{Der}(s(\num n)\doteq s(\num n)).
    \]
    For condition b), we have:
    \[
        \varphi[\num{n}/x]\equiv y\doteq s(\num n).
    \]
    Since $\num{\fun{succ}(n)}\equiv s(\num n)$, the implication reduces to:
    \[
        y\doteq s(\num n)\imp y\doteq s(\num n).
    \]
    Consequently,
    {\small
    \begin{align*}
        d_1&\colon\fun{Der}\bigl(
            y\doteq s(\num n)\imp y\doteq s(\num n)
        \bigr)
            &&;\ \text{Lem.~\ref{lem:identity}}\\
        d_2&\colon\fun{Der}\bigl(
            \forall y\,(
                y\doteq s(\num{n})\imp y\doteq s(\num n)
            )
        \bigr)
            &&;\ \fun{gen}(d_1).
    \end{align*}}
\end{xmpl}

\subsection{Primitive Recursive Functions}

To rigorously establish that metatheoretic decidable relations such as the proof predicate can be internalized, we formalize the class of primitive recursive functions as an inductive family of types $\fun{PR}(k)$ indexed by their arity $k\colon\N$. Elements of this type code for functions $\N^k\to\N$.

\begin{defn}[Primitive Recursion]
 The family $\fun{PR}\colon\N\to\univ$ is defined inductively by the following constructors:
 \begin{itemize}
 \item \textsl{Zero}: $\fun{zero}\colon\fun{PR}(0)$.
 \item \textsl{Successor}: $\fun{succ}\colon\fun{PR}(1)$.
 \item \textsl{Projections}: For any $k\colon\N$ and $i\colon\fun{Fin}(k)$, $\fun{proj}_i^k\colon\fun{PR}(k)$.
 \item \textsl{Composition}: For any $k, m\colon\N$, given $g\colon\fun{PR}(m)$ and a finite sequence of functions $h\colon\fun{Seq}(\fun{PR}(k), m)$, the composition is $\fun{comp}_{k,m}(g, h)\colon\fun{PR}(k)$.
 \item \textsl{Primitive Recursion}: For any $k\colon\N$, given $g\colon\fun{PR}(k)$ and $h\colon\fun{PR}(k+2)$, the recursion operator yields $\fun{prec}_k(g, h)\colon\fun{PR}(k+1)$.
 \end{itemize}
\end{defn}

\begin{rem}
 Each constructor in $\fun{PR}(k)$ determines a metatheoretic evaluation map $\fun{eval}_k\colon\fun{PR}(k)\to(\N^k\to\N)$ by structural recursion. Proving that all primitive recursive functions are representable in Peano Arithmetic requires demonstrating that for every $k\colon\N$ and $f\colon\fun{PR}(k)$, the function $\fun{eval}_k(f)$ satisfies the representability criteria.
\end{rem}

\subsection{Internalizing the Proof Predicate}

The culmination of syntactic representability allows the metatheory to reflect its own proof structures inward. 

\begin{thm}[Representability of the Proof Predicate]
 Because the evaluation maps of all primitive recursive functions are representable, and because the proof predicate\/ $\fun{Proof}(n,m)$ is constructed via bounded minimalization and primitive recursion, there exists a specific formula\/ $\fun{Prf}\colon\fun{Form}$ with two free variables\/ $x,y\colon\fun{Var}$ that represents it. Thus, for any natural numbers\/ $n,m\colon\N$, the following implication holds:
 \[
 \fun{Proof}(n,m) \to \fun{Der}(\fun{Prf}[\num{n}/x][\num{m}/y]).
 \]
\end{thm}